\section{吃螃蟹的人：从清华实验室到科技成果转化}

访谈开头，张鹏回到智谱的起点。智谱不是凭空创业，而是从清华大学计算机系知识工程实验室的研究和工程转化传统中长出来的。上一代人工智能四小龙让实验室看到 AI 产业化的可能，但也让他们意识到：感知智能阶段的能力天花板明显，下一代人工智能需要新的方向。

\begin{figure}[H]
\centering
\includegraphics[width=0.96\textwidth]{zhipu-ai-roadmap.png}
\caption{智谱六年路线：从认知智能、P2P 转化到上市节点，再回到 AGI 目标。}
\end{figure}

\begin{knowledgebox}{读图：这不是一条纯技术路线}
图中每个节点都同时包含技术和组织含义。认知智能是技术方向；P2P 是成果转化机制；GPT-3 和 ChatGPT 是外部范式冲击；DeepSeek 改变开源与效率叙事；IPO 则把公司带入新的治理阶段。
\end{knowledgebox}

\subsection{P2P：Paper to Product}

前面讲的是智谱的宏观路线，这里进入它最底层的组织方法。P2P 解释了为什么智谱天然会把论文、系统、产品和客户放在同一条链上。

张鹏反复讲到清华的 P2P 传统：paper to project / paper to product。研究成果不能只停在论文和 prototype，而要转成企业客户可用的 system 和 product。他在实验室里负责的正是工程转化：教授和学生做研究，工程团队把研究成果承接成系统，再交付客户需求。

\begin{figure}[H]
\centering
\includegraphics[width=0.94\textwidth]{paper-to-product-loop.png}
\caption{P2P：从论文到产品，再通过客户反馈反哺研究。}
\end{figure}

\begin{knowledgebox}{读图：P2P 是智谱的底层组织基因}
图中从 Paper 到 Prototype，再到 System、Product 和 Feedback。它说明智谱不是纯学院派研究院，也不是纯销售型公司；它的早期能力来自把前沿研究工程化，再让真实客户的需求反过来修正研究。
\end{knowledgebox}

\subsection{科技成果转化的第一只螃蟹}

前面讲 P2P 是能力，本节进一步看制度路径。科技成果转化如果没有走通，前沿研究即使能做成系统，也很难变成可以长期融资、经营和上市的公司资产。

2018 年左右国家政策为科研院所人员科技成果转化打开窗口，但细节并不清楚：成果如何估值，学校和团队如何分配，流程如何合规，谁来承担后续责任。智谱团队成了“吃螃蟹的人”：一边和学校谈，一边摸索制度路径，一年多后才把公司注册和团队转出走通。

\begin{warningbox}{成果转化不是简单“老师开公司”}
高校成果转化涉及知识产权、国资、学校权益、创始团队激励和公司治理。如果制度路径不清楚，后续融资、上市和合规都会留下隐患。智谱选择走正式路径，是公司长期化的重要前提。
\end{warningbox}

\subsection{制度路径为什么影响技术公司}

技术公司常常把注意力放在模型、论文、融资和产品上，但智谱的早期故事提醒我们：制度路径本身也是能力。如果一家从高校或科研院所转化出来的公司，知识产权归属不清、股权分配不清、学校权益不清，短期可能看不出问题，长期到了融资、审计、上市和国际合作时，就会成为结构性风险。

智谱花一年多时间做这个事，看起来慢，但它把未来风险提前消化掉了。这个选择和公司后来的“水泥感”一致：不追求最快的声量，而是把路基铺稳。对 AI 公司来说，这种低调的制度工作不如模型发布吸睛，却决定公司能不能走远。

\begin{knowledgebox}{科技成果转化的四个问题}
\begin{tabularx}{\textwidth}{p{0.22\textwidth}p{0.34\textwidth}X}
\toprule
问题 & 要解决什么 & 如果没解决会怎样 \\
\midrule
成果归属 & 论文、代码、专利和系统属于谁 & 融资和上市时产生产权争议。 \\
估值定价 & 科研成果折算成多少商业价值 & 股权分配和国资合规难以解释。 \\
团队激励 & 创始团队和学校如何分配权益 & 早期贡献者动力不足或后续纠纷。 \\
合规路径 & 是否符合学校和政策要求 & 发展越大，历史问题越难补。 \\
\bottomrule
\end{tabularx}
\end{knowledgebox}

\subsection{本章小结}

智谱的早期并不是“模型创业”，而是研究工程化和科技成果转化。P2P 基因解释了它为什么一开始就重视研究、系统、客户和商业化之间的闭环。

